\chapter{Functional Programming}
\label{ch:functional-programming}

\textit{Lambdas, \texttt{std::function}, \texttt{std::bind}, and reference wrappers --- the tools that let C++ treat behavior as a first-class value.}

C++11 brought functional-style programming into the mainstream. \textbf{Lambda expressions} create function objects inline; \textbf{\texttt{std::function}} stores any callable behind a uniform type; \textbf{\texttt{std::bind}} performs partial application; and \textbf{\texttt{std::reference\_wrapper}} lets references travel through these value-semantic facilities.

%% ============================================================
\section{Lambda Expressions}

A \textbf{lambda expression} creates an anonymous function object inline. Formally it is a prvalue whose result is a \textbf{closure object} behaving like a functor. The general syntax is \texttt{[captures](params) -> return\_type \{ body \}}.

\begin{lstlisting}[language=C++, caption={a basic lambda}]
auto add = [](int a, int b) { return a + b; };
int result = add(2, 3);   // 5
\end{lstlisting}

\begin{lstlisting}[language=C++, caption={a lambda passed to an algorithm}]
std::vector<int> v = {5, 2, 8, 1};
std::sort(v.begin(), v.end(), [](int a, int b){ return a > b; }); // descending
\end{lstlisting}

%% ============================================================
\section{Capture Lists}

By default a lambda cannot access enclosing-scope variables. The \textbf{capture list} makes selected variables part of the closure, by copy or by reference.

\begin{center}
\begin{tabular}{|l|l|}
\hline \textbf{Capture} & \textbf{Meaning} \\ \hline
\texttt{[]}      & Capture nothing \\
\texttt{[x]}     & Capture \texttt{x} by value \\
\texttt{[\&x]}   & Capture \texttt{x} by reference \\
\texttt{[=]}     & Capture all used variables by value \\
\texttt{[\&]}    & Capture all used variables by reference \\
\texttt{[this]}  & Capture the enclosing object \\
\hline
\end{tabular}
\end{center}

\begin{lstlisting}[language=C++, caption={capture by value vs reference}]
int a = 0;
auto f1 = [ ]() { return a * 9; }; // ERROR: 'a' not captured
auto f2 = [a]() { return a * 9; }; // OK: captured by value
auto f3 = [&a]() { return a++; };  // OK: by reference (modifies original)
\end{lstlisting}

\begin{quote}
\textbf{Dangling-capture warning:} a reference (or \texttt{this}) capture must not outlive the captured object. A \texttt{[\&]} lambda stored in a \texttt{std::function}, thread, or container is a classic use-after-free. Prefer value capture for anything that may outlive the frame.
\end{quote}

%% ============================================================
\section{Mutable Lambdas}

Value captures are \texttt{const} by default; \texttt{mutable} removes that, letting the lambda modify its own copy (state retained across calls of the same closure, never touching the original).

\begin{lstlisting}[language=C++, caption={a mutable lambda holds state}]
int x = 0;
auto increment = [x]() mutable { return ++x; };
increment(); // 1
increment(); // 2  (state retained); outer x still 0
\end{lstlisting}

%% ============================================================
\section{The Full Lambda Anatomy}

\begin{center}
\begin{tabular}{|l|l|}
\hline \textbf{Part} & \textbf{Role} \\ \hline
default-capture     & \texttt{=} or \texttt{\&}; precedes the list \\
capture-list        & per-variable capture \\
argument-list       & the parameters \\
\texttt{mutable}    & makes value captures non-\texttt{const} \\
throw-spec          & e.g. \texttt{noexcept} \\
attributes          & e.g. \texttt{[[noreturn]]} \\
\texttt{-> type}    & explicit return type (when not deducible) \\
lambda-body         & the implementation \\
\hline
\end{tabular}
\end{center}

\begin{lstlisting}[language=C++, caption={explicit return type and noexcept}]
auto divide = [](double a, double b) noexcept -> double {
    return b != 0.0 ? a / b : 0.0;
};
\end{lstlisting}

Each lambda has a unique, unnamable type --- which is why you store them in \texttt{auto}, templates, or \texttt{std::function}.

%% ============================================================
\section{\texttt{std::function}: The Polymorphic Callable Wrapper}

\texttt{std::function<R(Args...)>} is a type-erased wrapper holding any callable matching its signature: free function, function pointer, lambda, functor, or \texttt{bind} result.

\begin{lstlisting}[language=C++, caption={one type, many callables}]
using fn_t = std::function<double(int, float, double)>;
double foo_fn(int x, float y, double z) { return x + y + z; }
struct Foo { double operator()(int x, float y, double z){ return x+y+z; } };

fn_t a = foo_fn;        // free function
fn_t b = Foo{};         // functor
fn_t c = [](int x, float y, double z){ return x+y+z; }; // lambda
std::vector<fn_t> table = {a, b, c};
for (auto& f : table) f(1, 2, 3);
\end{lstlisting}

%% ============================================================
\section{\texttt{std::bind} and Partial Application}

\texttt{std::bind} fixes some arguments of a callable; unbound positions use placeholders \texttt{\_1}, \texttt{\_2}, \ldots

\begin{lstlisting}[language=C++, caption={partial application}]
using namespace std::placeholders;
int add(int a, int b) { return a + b; }
auto add5 = std::bind(add, 5, _1); // fix first arg
int r = add5(10);                  // add(5,10) -> 15
\end{lstlisting}

\begin{lstlisting}[language=C++, caption={binding a member function and reordering}]
struct Calc { double weighted(int x, double z, float y){ return x+y+z; } };
Calc c;
auto g = std::bind(&Calc::weighted, &c, _1, _3, _2);
g(1, 2.0f, 3.0); // calls c.weighted(1, 3.0, 2.0f)
\end{lstlisting}

\begin{quote}
\textbf{Modern guidance:} lambdas largely replace \texttt{std::bind} --- more readable, easier to inline. Use \texttt{bind} only when a lambda would be clumsier.
\end{quote}

%% ============================================================
\section{\texttt{std::reference\_wrapper} and \texttt{std::ref}}

Value-semantic facilities copy their arguments. To send a \emph{reference} through them, wrap it with \texttt{std::reference\_wrapper<T>} via \texttt{std::ref} (or \texttt{std::cref} for const).

\begin{lstlisting}[language=C++, caption={forcing reference semantics through a copying interface}]
void increment(int& n) { ++n; }
int x = 0;
auto bound = std::bind(increment, std::ref(x)); // bind copies -> use ref()
bound();                                          // x is now 1
\end{lstlisting}

\begin{lstlisting}[language=C++, caption={a container of references}]
int a = 1, b = 2, c = 3;
std::vector<std::reference_wrapper<int>> refs = {a, b, c};
for (int& r : refs) r *= 10;   // a,b,c -> 10,20,30
\end{lstlisting}

A \texttt{reference\_wrapper} is copyable/assignable, implicitly converts back to \texttt{T\&}, and is how you store references in a \texttt{vector} or pass references to \texttt{std::thread}.

%% ============================================================
\section{Callback Patterns}

\begin{lstlisting}[language=C++, caption={object B registers a member-function callback on object A}]
class A {
public:
    std::function<void(int, const std::string&)> on_event;
    void fire() { if (on_event) on_event(100, "event fired"); }
};
class B {
    A a_;
public:
    B() { a_.on_event = [this](int i, const std::string& s){ handle(i, s); }; }
    void handle(int i, const std::string& s) { /* ... */ }
    void run() { a_.fire(); }
};
\end{lstlisting}

%% ============================================================
\section{Professional Insights}

\textbf{\texttt{std::function} is not free.} Its value semantics force a copy/move of the callable, frequently with a heap allocation. Small-object optimization is implementation-defined and only for \texttt{noexcept}-movable types. The call is indirect --- about a virtual-call cost.

\textbf{Prefer a template parameter in hot paths.} If you only invoke a callable (e.g. a comparator), take it as a template parameter so it is monomorphized and inlined; reserve \texttt{std::function} for genuine type erasure (callback slots, command tables, plugin hooks).

\textbf{Lambdas inline; \texttt{bind}/\texttt{function} often do not.} Keep callables concrete up to the boundary where erasure is genuinely required.

\textbf{Capture deliberately.} \texttt{[=]} can silently copy expensive objects; \texttt{[\&]} can dangle. Name captures to document intent and prevent both bugs.
